% ED26-BM-RAT2021-intro: source-based treatment of sections 1--9.
\chapter{Frameworks for Assessing Rationality}
\label{sec:p11-rationality-intro}

Rationality is an object of philosophical assessment and psychological
investigation, not one universally interchangeable score. Knauff and Spohn's
2021 systematic introduction asks how these perspectives can inform each
other while preserving their different questions
\cite[pp.~1--10]{KnauffSpohn2021FrameworksIntro}.

\section{Why rationality is worth investigating}

The authors begin with the conception of humans as rational animals,
then consider three objections rather than treating that conception as
an empirical finding. First, apparently irrational conflict, prejudice
and destructive collective behavior challenge any identification of
rationality with universal successful performance. Their response
distinguishes possessing a capacity from exercising it successfully on
every occasion. The judgment of an assessor and the level at which
an activity is assessed also matter.

Second, rationality-related capacities are not confined to humans.
The authors recognize basic forms in other species and evolutionary
continuities while reserving a fuller human conception for richer
representational and reflective abilities. This is their framework
position, not a supplied comparative experiment that measures every
species against one invariant scale.

Third, language, art, religion, education and other capacities might
seem equally suitable to characterize humans. Their language example
argues that even its varied social purposes depend on the formation
and exchange of beliefs and desires, so a rationality-related dimension
cannot simply be removed from understanding it. This is not a claim
that every utterance is sincere, every linguistic activity successful,
or all conduct consciously reasoned.

Their motivating examples of war, starvation, climate change, pandemics
and technological difficulties are reasons to examine how decisions
and critical reflection work, not estimates of the relative amount of
rationality in a population. The authors expressly deny that rationality
alone suffices: empathy, justice, honesty, appropriate values and other
qualities can matter. Their further claim is that proposed means and
ends must still be open to reasoned assessment. These are philosophical
and practical commitments of the introduction, not a demonstrated
formula making rationality identical with every virtue
\cite[pp.~1--2]{KnauffSpohn2021FrameworksIntro}.

\begin{hbdefinition}{def:p11-ri-capacity-performance}{Capacity, performance and assessment target}
A capacity or competence is an ability to produce an appropriate
performance under relevant conditions. A performance is what occurs
on a particular occasion. Inferring that a capacity is absent from
one departure requires an account of the task, interpretation and
conditions; inferring that performance is always appropriate from
the possession of a capacity requires another argument.

An assessment also identifies its target: an individual belief or
action, a process, a joint result, or a disposition. Switching target
can change the question rather than merely improve the same score.
The source invokes Chomsky's competence/performance distinction in
this initial clarification; it does not provide one operational measure
of rational competence for all tasks.
\end{hbdefinition}

\begin{hbdefinition}{def:p11-ri-two-agent-model}{A stipulated individual/joint contrast}
Two players choose \(\mathrm{Coop}\) or \(\mathrm{Def}\).
The row gives player 1's action, the column player 2's, and each pair
is their respective stipulated payoff:
\[
\begin{array}{c|cc}
 &\mathrm{Coop}&\mathrm{Def}\\\hline
\mathrm{Coop}&(3,3)&(0,4)\\
\mathrm{Def}&(4,0)&(1,1)
\end{array}
\]
The individual criterion is to select a best response to the other's
action, meaning that no unilateral switch increases one's own payoff.
The joint comparison used here asks whether both payoffs can be increased
at another profile. Neither criterion is declared a universal definition
of rationality, and the table is not participant data.
\end{hbdefinition}

\begin{hbtheorem}{thm:p11-ri-individual-joint}{proved=annex:p11.ri.individual-joint}{Individual best responses need not produce a jointly beneficial result}
In this model \(\mathrm{Def}\) is strictly better for each player
against either action of the other. The unique mutual-best-response
profile is \((\mathrm{Def},\mathrm{Def})\), with payoff \((1,1)\).
Both players would receive a higher payoff at
\((\mathrm{Coop},\mathrm{Coop})\), namely \((3,3)\).
Thus conformity to the stated individual criterion does not entail
the stated joint benefit. This construction is not evidence that
typical people defect or that every collective disaster is caused
by individually rational choice.
\end{hbtheorem}

The limited result explains why identifying a disaster does not by
itself locate the defect in each person's inference. Incentives,
coordination, institutional choices and the joint criterion would need
to be specified. The introduction's groupthink and disaster references
are broader literature illustrations; this table neither reconstructs
those events nor measures their causes.

\section{The interdisciplinary commitment}

The introduction rejects a simple history in which philosophy alone
addressed rationality for millennia and psychology suddenly began at
the end of the nineteenth century. Earlier philosophical psychology
also addressed mind and cognition, with less clearly separated
methodologies. The authors contrast psychology's later empirical
specialization with philosophy's clearer normative focus and describe
the gap that their interdisciplinary project seeks to bridge.

They discuss two contrary research positions. One permits descriptive
psychology to proceed independently of normative work. Another appeals
to competence and performance to argue that experimental departures
cannot refute human rationality. The editors instead commit to combining
empirically assessed descriptive theories with elaborated normative
standards as points of comparison. Their description of these positions'
standing in 2021 is not a measured contemporary consensus.

Describing how people respond, specifying what an appropriate response
would be, and justifying the authority of that standard are different
activities. The actual Broome treatment distinguishes correctness from
having a normative reason to conform, and motivating from normative
reasons (see \cref{def:p11-reasons-normative-senses,def:p11-reasons-roles}).
The empirical benchmark/response distinction makes a related separation
(see \cref{def:p19-benchmark}). Neither interface requires equating
an observed regularity with an obligation or assuming that any stipulated
benchmark has already been justified
\cite[p.~2 and note~1]{KnauffSpohn2021FrameworksIntro}.

\section{Contours before a universal definition}

The source begins its conceptual account with a broad field: acquiring
concepts, belief formation, knowledge, inference, judgment, planning,
choice, calculation and the satisfaction of needs or desires. It does
not infer that one formal model already covers that whole field.
Its discussion of Greek, Latin, French, English and German terminology
warns that inherited terms have changing and doctrine-dependent meanings.
Kant's distinction between understanding and reason is not simply the
same contrast as hermeneutic understanding and explanation.

\begin{hbdefinition}{def:p11-ri-reason-senses}{Faculty, consideration and activity}
The source separates \emph{reason} as a cognitive faculty,
\emph{a reason} as a consideration offered to explain or justify
something, and \emph{reasoning} as an activity of inferring or
arriving at a conclusion. Giving reasons need not be the same
activity as presently reasoning; a hypothetical inference can also
proceed from premises whose status is not endorsed.

In the introduction's historical comparison these English senses
correspond to different German terms, including \emph{Vernunft},
\emph{Grund} and expressions for reasoning or drawing a conclusion.
The reported Greek/Latin history, Cicero's use of \emph{ratio} for
\emph{logos}, the Greek rational-number connection, and Tertullian's
early use of \emph{rationalitas} are conceptual
orientation in the source, not an independently undertaken etymological
investigation. It reports that the faculty terms remained more prominent
than \emph{rationality} until that latter term's increased use in
nineteenth-century economics, social science and normative discussion.
Its note 4 points to Broome's related treatment.
\cite[pp.~3--4 and notes~3--4]{KnauffSpohn2021FrameworksIntro}
\end{hbdefinition}

The faculty sense therefore does not establish that every explaining
consideration is a normative reason or that rationality reduces to
correct response to objective reasons. Broome's specifically structural
conception and mental-base qualifications remain distinct
\cref{def:p11-reasons-structural}. The introduction is an organizing
perspective alongside such accounts, not their replacement.

\section{Historical milestones and what they compare}

The introduction selects milestones rather than giving a complete
history. Boole's laws of thought concern mental operations, whereas
Frege's antipsychologism separates psychological processes from logical
truth and normativity. Frege's use of \emph{thought} refers to objective
content, not simply an occurrence in someone's mind. The source contrasts
the normative prescription associated with logical laws and Wundt's
psychological-function account. It places Frege's 1879 project beside
Wundt's Leipzig institution in the same year.

Its Störring illustration concerns a 1908 deductive-problem battery,
with verbal responses, response times, eye movements, gestures and
respiration among recorded observations. The introduction supplies no
study sample, recruitment, detailed scoring or uncertainty for those
results. The example locates a descriptive research program; it is
not a newly qualified original experiment here.

Piaget's genetic epistemology is presented through accommodation and
assimilation, with sensorimotor, preoperational, concrete-operational
and formal-operational stages. The source's age ranges are birth to
2, 2--7, 7--11 and 11 onward. These are reported features of that
developmental theory, not an assertion that every child's actual
trajectory has those boundaries. Its rule-based account of later
reasoning is contrasted with Johnson-Laird's mental-model approach
and their different proposed representations and processes.

The Popper--Kuhn discussion makes the benchmark question consequential.
The source contrasts falsificationist scientific methodology with
observations of actual scientific practice, including the controversies
with logical positivism and the Frankfurt school. One response to
departures is to judge practice irrational; another questions the
adequacy of the proposed norm. Lakatos's research-program account is
presented as one attempted response. Neither actual prevalence nor a
scientist's reputation alone settles which response is justified.
The introduction relates this debate to defeasible reasoning and
theory change without treating every cited scientific methodology
as fully reconstructed here.

\begin{hbdefinition}{def:p11-ri-card-task}{The actual selection-task illustration}
Each of four displayed cards has a letter on one side and a number
on the other. Their visible faces are \(A,D,4,7\). The proposed
rule says that a card with a vowel has an even number on its other
face. Let \(\mathcal C\) be these four card objects,
\(\mathsf V(k)\) mean its letter is a vowel, and
\(\mathsf E(k)\) mean its number is even. For the closed task,
the rule is
\[
 \forall k\in\mathcal C\,(\mathsf V(k)\to\mathsf E(k)).
\]
The intended instruction is to expose any possible violation among
the displayed cards, with the two-side setup taken as given. It is
not a sampling model for an unobserved population or a deontic
permission rule.
\cite[p.~5]{KnauffSpohn2021FrameworksIntro}
\end{hbdefinition}

\begin{hbtheorem}{thm:p11-ri-card-check}{proved=annex:p11.ri.card-check}{Which hidden faces can violate the stated rule}
The only violating pattern is vowel with odd number. For these four
cards, turning the \(A\) and the \(7\) cards is sufficient
to check every still-possible violation. Omitting either leaves
a possible unchecked violation; the hidden faces of \(D\) and
\(4\) cannot violate the rule under the stipulated setup.
The conclusion concerns this closed task and this conditional
interpretation, not every conversational use of an if-clause.
\end{hbtheorem}

The introduction reports a low success frequency from the selection-task
literature. Its expression in terms of fewer than ten out of a hundred
is not an identified experiment with \(N=100\): no particular sample,
recruitment, estimate or confidence interval is supplied there. The
normative card calculation and a population claim therefore have
different evidence requirements. An incorrect final selection also
does not identify which process, prior assumption or interpretation
produced it. The current evidence record makes this distinction explicit
\cref{def:p19-empirical-question,def:p19-benchmark}.

The source next describes belief-bias paradigms, in which validity and
believability are varied separately, and the later emphasis on subjective
degrees of belief. What counts as a prior-belief bias in one task can
be relevant information in another. Its description of a new Bayesian
paradigm is accompanied by qualifications: note 6 questions a literal
Kuhnian paradigm-shift interpretation, and note 7 observes that
Bayesian labels have different uses. The milestone is a change of
research questions and proposed norms, not a universal diagnosis of
every nonclassical response.

The probability strand includes the updating idea associated with Bayes,
with the source's priority qualification, Ramsey's 1926 subjective
probability and utility program, and de Finetti's contribution.
The finite Bayes theorem is a specified mathematical result
\cref{thm:p15-finite-bayes}; it does not show that every person's
uncertainty is probabilistic, or that merely eliciting numerical confidence
supplies all the model's assumptions. The introduction explicitly allows
other representations and revisions of graded belief.

Simon's bounded-rationality program brings limited knowledge, memory,
representation and computation into the question. Satisficing stops
search when an option reaches an aspiration level; it is not thereby
a proof that the first such option globally maximizes a known utility.
The introduction contrasts heuristics-and-biases, ecological rationality
and prospect theory, with their different descriptive and normative
emphases. Their reported empirical advantages are literature-level
comparisons here, not a supplied common-task meta-analysis.
The current bounded-rationality treatment likewise distinguishes a
principle of limits, positive mechanisms and prescriptions
\cref{def:p17-br-principle-theory,def:p17-br-three-rationalities}.

\section{What is being assessed, and relative to what?}

\begin{hbdefinition}{def:p11-ri-assessment-objects}{Primary, derived and arational assessments}
The source distinguishes objects assessable as rational or irrational
from \emph{arational} phenomena outside that assessment, such as
the weather. Its primary focus is beliefs and other epistemic attitudes,
and actions; rational descriptions of a road design or supermarket
layout are derived through their connection with such activities.
The assessment of emotions is a further disputed question.

Here \emph{epistemic} has the source's deliberately broad doxastic
sense, including opinions, doubts and credences. Note 8 distinguishes
that use from a strict knowledge-specific sense; certainty, true
belief and knowledge are not made identical by this terminology.
The source contrasts controlled or intentional action with mere
behavior: an involuntary sneeze is not assessed in the same way as
an intentional imitation. Habitual activity need not involve a
consciously entertained intention on each occasion. Intention and
control are therefore delicate conditions, not a complete sufficient
definition of action.
\cite[pp.~7--8 and notes~8--9]{KnauffSpohn2021FrameworksIntro}
\end{hbdefinition}

The umbrella example distinguishes a categorical assessment from an
assessment relative to the person's information. A rain belief can
look unreasonable in light of current evidence but intelligible relative
to an outdated forecast; an umbrella choice can follow reasonably from
that belief without the belief itself being well supported. Specifying
the background does not automatically establish its adequacy.

The same issue occurs when an experiment instructs a participant to
take its premises as fixed facts. A response may violate that benchmark
because the person reasons incorrectly under it, or because the person
imports another premise or interpretation. That possibility does not
prove that the response is rational, but it prevents diagnosing the
process from a verdict alone. Defeasible revision under added information
is not a counterexample to the monotonicity of a fixed classical
semantic consequence relation.

The authors treat relative assessments as basic and categorical
assessments as depending additionally on the standing of their background.
In their epistemic discussion, beliefs are supported by information
and other epistemic attitudes; a wish can cause a belief without
constituting evidence for it. They identify wishful thinking as the
relevant problem, not as an algorithmically established classification
of every emotion-influenced belief. Note 11 discusses mediation by
perception and belief even in fact-based accounts of reasons; note 13
connects the further regress question with the Agrippa trilemma and
its competing philosophical responses. Those responses are not
settled by the relative/categorical terminology.

For actions, instrumental rationality assesses how a choice serves
goals in light of beliefs and pro-attitudes. \emph{Desire} is used
as a broad umbrella for wishes, interests, goals, norms and related
driving attitudes, not a reduction to one psychological variable.
The source distinguishes stable dispositions from situation-specific
forces and then intrinsic ends from extrinsic wants serving other ends.
Its aspiring-millionaire/philosopher case exposes the need to justify
a means belief or the consistency of an extrinsic aim with its end;
it is not income data or an unconditional career recommendation.

Whether intrinsic desires themselves admit rational assessment is
contested. The introduction contrasts Hume's instrumental position
with Kant's connection of practical reason and morality, Weber's
value/instrumental contrast and Habermas's further social conception;
it also reports a contrary instrumental reading of communication.
This is not a claim that those theories share one concept of rationality.
The authors recognize the difficulty of separating rational and moral
assessment and explicitly limit the depth of their treatment.
The supplied ordering in \cref{def:p17-weak-preference} is likewise
not an independent justification of a person's ultimate ends.

\section{Levels of explanation and matching comparison interfaces}

The source closes its preliminary concepts with Marr's distinction:
what a system is to achieve, the representations and procedure by
which it achieves it, and physical implementation. It discusses
biological naturalism, functionalist indifference to the implementation
medium, and views resisting even functional reduction as different
positions, not one settled account of mind.

The authors adopt \emph{output-oriented} and \emph{process-oriented}
for the first two levels, with a deliberate shift that permits a
normative specification of the goal. These are not merely alternate
words for running a computer program. A correct output, a sound
procedure and an actual implementing mechanism answer different
questions. This distinction is developed below in
\cref{sec:p11-ri-process-output}; it is not exhausted by these
preliminary terms.

The current attitude-based distinction between theoretical and
practical rationality supplies a matching orientation
\cref{def:p11-rat22-attitude-types}, while Broome's mental-base
qualification reminds the reader to say which attitudes and worldly
relations are held fixed. The psychological evidence treatment asks
which construct the task and measurement support. Linking these
approaches requires their actual assumptions, not a new generic
taxonomy that silently merges them.

\begin{hbmethod}{def:p11-ri-opening-assessment}{Specify the assessment before applying a standard}
State the target and occasion, the information and pro-attitudes in
play, and whether the verdict is relative or all-things-considered.
Distinguish possession of a capacity from the recorded performance;
separate what occurred from the benchmark and its justification.
For a formal result, state the model and its boundary. For a
historical or empirical claim, retain its actual source, task,
sample and uncertainty, marking unavailable details as unavailable.
Do not infer a whole population's rationality, or an agent's complete
process, from one answer or one mathematical representation.
\end{hbmethod}

\section{Theoretical and practical rationality}
\label{sec:p11-ri-theoretical-practical}

The introduction separates two questions with a demonstration example.
First one assesses the claim that a herbicide harms bees; then, taking
one's resulting belief into account, one decides whether to attend a
protest. The first question calls for evidence and inference. The second
also involves motives: concern for nature, an anticipated opportunity
to meet friends, the attraction of a quiet Sunday, and possible weather
or safety costs. The example stipulates a conclusion for the purpose
of continuing deliberation; it does not establish a toxicological result
or recommend participation in a particular event
\cite[pp.~10--11]{KnauffSpohn2021FrameworksIntro}.

\begin{hbdefinition}{def:p11-ri-direction-fit}{Attitude type and direction of fit}
Theoretical rationality concerns the justification of beliefs and other
epistemic attitudes, their formation and revision. Practical rationality
concerns actions, intentions and pro-attitudes, and therefore requires
both epistemic and motivational considerations. In the source's
direction-of-fit terminology, belief aims to represent how the world
is, whereas action aims to make the world accord with a purpose.
Reasoning can occur on both sides; its occurrence does not erase the
difference between the attitudes and their justifications.
\end{hbdefinition}

The authors acknowledge approaches that integrate reasoning with
decision making through utility and social pragmatics. They also discuss
doxastic voluntarism and epistemic decision theory, in which belief
formation is understood as choice serving epistemic aims. Their own
position is nevertheless asymmetric: practical rationality presupposes
theoretical rationality, since action assessment includes beliefs about
the situation and the effectiveness of means, while epistemic assessment
need not include the agent's practical motives. This is the introduction's
philosophical position, not a demonstrated separation of two independent
brain systems. Wedgwood's attitude-based distinction provides a related
comparison without making the subject matter of a proposition determine
its attitude type (\cref{def:p11-rat22-attitude-types}).

\subsection{Forms and changes of epistemic attitudes}

The source takes beliefs to have contents, conventionally called
propositions, rather than identifying a belief with a particular sentence
token: speakers of different languages can share a belief. It deliberately
does not settle the disputed ontology of propositions or concepts.
Its note~16 preserves that restriction. An epistemic state can then be
described qualitatively, comparatively or quantitatively
\cite[pp.~11--14 and notes~16--18]{KnauffSpohn2021FrameworksIntro}.

\begin{hbdefinition}{def:p11-ri-epistemic-formats}{Three formats, not three interchangeable measures}
A qualitative description records accepted contents. The source calls
an unrestricted collection a \emph{belief base}; a consistent,
deductively closed collection is a \emph{belief set}. The additional
conditions are substantive rationality postulates, not facts about every
person's stored beliefs.

A comparative description orders contents by credibility or plausibility.
A quantitative description assigns degrees. A numerical assignment is
not a probability merely because its values lie between zero and one:
probability axioms impose further conditions. The source's Bayesian
format uses subjective probabilities, while recognizing alternative
formats and disputing the claim that this one format exhausts epistemic
attitudes.
\end{hbdefinition}

The Lockean threshold proposal connects acceptance with a degree above
a threshold. The authors report the controversy rather than adopting
that reduction. The following finite construction makes one obstruction
precise. It does not refute every context-sensitive or non-threshold
account of belief, and its events are stipulated, not elicited beliefs.

\begin{hbdefinition}{def:p11-ri-threshold-model}{A fixed-threshold acceptance model}
Fix \(0\leq t<1\). On a finite uniform probability space
\(\Omega=\{1,\ldots,n\}\), accept an event \(E\) exactly when
\(\mathbb P(E)>t\). Let \(E_i=\Omega\setminus\{i\}\).
Closure under finite conjunction would require acceptance of
\(\bigcap_{i=1}^n E_i\) whenever every \(E_i\) is accepted.
\end{hbdefinition}

\begin{hbtheorem}{thm:p11-ri-threshold-closure}{proved=annex:p11.ri.threshold-closure}{A fixed threshold below certainty does not ensure conjunctive closure}
For each such \(t\), there is an integer \(n\geq2\) for which
every \(E_i\) is accepted but their conjunction is not. Thus this
threshold rule on all finite uniform spaces cannot by itself produce
the source's deductively closed belief sets.
\end{hbtheorem}

Static form is only one question. The introduction distinguishes
\emph{internally driven} changes through thinking from
\emph{externally driven} changes through perception, testimony or new
information. On the process side, an account must explain the
representations and operations that transform the input, including
possible errors. On the output side, one asks which conclusion or
revised state is appropriate under the stated input. A sound formal
procedure is not, solely by being sound, a psychologically observed
procedure; different procedures may realize the same input--output
relation. Nonmonotonic and defeasible procedures, as well as deductive
ones, belong to this comparison.

Externally driven revision raises questions about conflict with prior
beliefs, information order, minimal change, entrenchment and source
reliability. The introduction relates the normative revision problem to
induction, while distinguishing it from a causal account of how a mind
processes the new input. Finite Bayes' formula supplies the conditional
probability identity when its hypotheses hold
(\cref{thm:p15-finite-bayes}); adopting conditionalization as the proper
revision rule requires a further normative and evidential argument.
Other epistemic formats have their own proposed learning rules.
Suppositional reasoning also complicates the internal/external division:
working out a supposition need not have the same consequences as learning
that it is actually true. The source identifies that issue without
settling all supposing-versus-updating theories.

\subsection{Motivational attitudes and practical change}

Practical assessment has parallel qualitative, comparative and
quantitative descriptions. One may record what is desired, order
alternatives by preference, or use a utility scale. The source groups
drives, wishes, goals, norms and values as pro-attitudes for exposition,
but leaves open whether they require different theories. This
terminological grouping neither makes moral obligations identical with
appetites nor proves that one scale measures every motive
\cite[pp.~14--16 and notes~19--21]{KnauffSpohn2021FrameworksIntro}.

Its content examples also caution against treating desires as simple
relations to extensional objects. Wanting to meet Jekyll but not Hyde
can involve different descriptions of one person; wanting a fictional
character need not identify an existing object. The source uses
propositional contents for convenience, not as a resolution of those
intentional-content problems.

The qualitative description must still account for the interaction
between intrinsic and extrinsic desires. Wanting to go without a coat
can depend on a belief about the weather, and that means preference
can change when the belief changes although the underlying end does
not. Comparative preference theory and numerical utility theory make
such interactions more tractable. Their representation results require
particular domains and axioms; an ordinal ranking of outcomes alone
does not determine preferences over risky mixtures. The current finite
preference and lottery treatments retain those requirements
(\cref{def:p17-weak-preference,thm:p17-finite-vnm}).

Expected-utility maximization combines probability and consequence
utility. It may specify which act should be selected without claiming
that the agent mentally multiplies and sums the displayed numbers.
The introduction contrasts that output reading with process-oriented
heuristic and bounded-rationality accounts. Empirical inadequacy of an
explicit-calculation model does not alone refute the same formula as
an output norm, just as numerical agreement alone does not identify
an actual calculation process.

Information-dependent intentions give a simple practical dynamic.
One can plan to wear a T-shirt if the forecast is warm and a coat
otherwise. The intention changes after information arrives, with the
goal of an appropriate outing held fixed. This differs from an
intrinsic-preference change caused by maturation, education, advertising,
satiety or boredom. The introduction places the latter in descriptive
motivation research while recognizing questions about morally proper
value development. It distinguishes such questions from dynamic
decision theory's question of how to act in view of changing preferences.
It does not claim that a single current-preference optimization settles
all conflicts between old and new standpoints.

\subsection{Truth-directed assessment without a truth guarantee}

The source's modified mother-and-son illustration separates an
epistemic assessment from a proposed moral or practical justification
of a belief. If evidence strongly indicates guilt, believing innocence
is not thereby epistemically supported because loyalty is valued.
The example's suggestion that maternal support might be morally
required is a conditional premise, not a moral theorem about parents.
The authors endorse the further normative position that nonepistemic
demands for epistemically irrational belief should be rejected
\cite[pp.~16--17 and notes~22--23]{KnauffSpohn2021FrameworksIntro}.

They give truth the central role in epistemic assessment while
acknowledging proposed further aims, including simplicity, relevance,
systematicity and explanatory power, and the alternative of taking
knowledge as the aim. Truth-directed rational formation is not
infallibility: misleading evidence can support a false conclusion,
while a guess can accidentally be true. Consequently the truth of
a conclusion does not establish the rationality of its formation,
and its falsity alone does not establish irrationality.

Credences also require a different evaluative form from outright
beliefs. A probability assignment is not itself a true or false
proposition about which event will occur; accuracy comparisons require
a specified scoring or accuracy account. Neither a binary truth verdict
nor the mere use of numbers settles that account. The introduction
does not resolve whether all epistemic aims reduce to truth.

The authors list correspondence, semantic, deflationary, pragmatic,
consensus, coherence and constructivist truth approaches as disputed
positions, not equivalent theories. Their comparison with psychological
constructivism distinguishes explaining how people take something to
be true from giving an account of truth itself. A constructive role for
experience and cognition does not, without another argument, identify
truth with belief. Likewise, the debate over moral realism and whether
morality is required by rationality prevents simply copying the
truth-directed aim onto every practical end. This is a boundary on
the present comparison, not an adjudication of all those truth theories.

\section{Normative and descriptive theories}
\label{sec:p11-ri-normative-descriptive}

The introduction asks three separate questions: what rationality
requires, how people actually reason and act, and how those inquiries
are related. Its distinction between empirical and genuine normativity
is important. Describing a community's accepted and sanctioned rule
records a social fact. Accepting a rule as a standard for one's own
conduct takes an internal normative stance. Hart's external/internal
contrast supplies the illustration; observed conformity alone does
not establish the rule's genuine authority
\cite[pp.~18--21]{KnauffSpohn2021FrameworksIntro}.

\begin{hbdefinition}{def:p11-ri-no-bridge-language}{An explicit no-bridge comparison}
Let \(\Gamma\) be a satisfiable set of propositional formulas built
only from atoms designated descriptive. Let \(N\) be a fresh atom
designated as a normative assertion, absent from those formulas.
Use ordinary classical valuations with no additional axiom linking
\(N\) to the descriptive atoms. Designating an atom records the
comparison's intended interpretation; it does not supply a deontic
logic or a theory of normative truth.
\end{hbdefinition}

\begin{hbtheorem}{thm:p11-ri-no-bridge}{proved=annex:p11.ri.no-bridge}{Independent assertions require a bridge premise}
Under this explicit model, \(\Gamma\) entails neither \(N\) nor
\(\neg N\). Conversely, if a descriptive atom \(D\) is distinct
from \(N\), neither \(N\) nor \(\neg N\) entails \(D\) or
\(\neg D\). Pure assertion of one side cannot determine the other
without constraints connecting them.
\end{hbtheorem}

This elementary result isolates a logical gap. It is not a theorem
that every mixed normative argument is invalid. Adding a substantive
bridge premise changes the premises, and inconsistent premises are
explicitly excluded. Note~30 distinguishes Hume's is--ought transition
from Moore's related question about goodness; the source's discussion
does not make those historical problems identical.

\subsection{Why the benchmark itself needs argument}

Classical deduction preserves truth from true premises. That feature
supports a proposed epistemic norm but neither guarantees true premises
nor obtains information not already implied by them. It also does not
by itself uniquely select classical logic over all competing accounts
of consequence. The introduction therefore regards truth preservation
as one consideration, not a completed universal norm for ampliative
or uncertain reasoning. Its discussion of conditionals, relevance,
nonmonotonic inference and AI illustrates the need to identify the
reasoning task rather than classifying every departure from material
implication as an error. Notes~24--27 preserve the qualification that
Gricean implicature provides a competing response; the historical claim
about logical empiricism's difficulties is the authors' interpretation.

Similarly, interpreting degrees as probabilities commits one to
probability laws, but that interpretation and its normative justification
are separate questions. The source recognizes accuracy arguments,
alternative degrees-of-belief formats and competing learning rules.
Probability axioms alone do not force inductive expectations from past
observations. For example, a stipulated joint law assigning probabilities
\(0.05,0.45,0.45,0.05\) respectively to
\((W_1,W_2),(W_1,\neg W_2),(\neg W_1,W_2),
(\neg W_1,\neg W_2)\) is coherent, yet gives
\(\mathbb P(W_2)=0.5\) and
\(\mathbb P(W_2\mid W_1)=0.1\). Here \(W_1,W_2\) are
two specified white-swan observation events. The model is deliberately
stipulated, not an empirical swan distribution; choosing an appropriate
law is exactly the additional issue.

The same normative dispute arises for decision theory. Representation
arguments connect particular preference conditions to expected utility;
the appropriateness of the conditions is not established by the
representation alone. The introduction's Ellsberg variant displays
the consequence of using a single additive prior and common
state-independent consequence utility
\cite[pp.~19--20]{KnauffSpohn2021FrameworksIntro}.

\begin{hbdefinition}{def:p11-ri-ellsberg-model}{The introduction's red, white and black bets}
An urn has 30 red balls and 60 black or white balls in an unknown
split. A uniformly drawn ball gives mutually exclusive exhaustive
events \(R,W,B\), with \(\mathbb P(R)=1/3\) and
\(\mathbb P(W)+\mathbb P(B)=2/3\). A bet on event \(E\)
gives the same reward \(H\) when \(E\) occurs and the same
otherwise consequence \(L\). Its expected utility is
\(u(L)+(u(H)-u(L))\mathbb P(E)\), where \(u(H)>u(L)\).
The pattern considered is strict preference for the \(R\) bet over
the \(W\) bet, but for the \(B\cup W\) bet over the
\(B\cup R\) bet. It is a pair of offered comparisons, not an
assumption that one ball can have two colors.
\end{hbdefinition}

\begin{hbtheorem}{thm:p11-ri-ellsberg-contradiction}{proved=annex:p11.ri.ellsberg-contradiction}{The common-event contradiction in the introduction's variant}
No single additive prior with the stated common consequence utility
represents both strict preferences. The first comparison requires
\(\mathbb P(R)>\mathbb P(W)\), the second
\(\mathbb P(W)>\mathbb P(R)\).
\end{hbtheorem}

The inference is a representation failure under specified assumptions,
not a diagnosis of every chooser's psychology. The source reports a
frequent preference pattern but supplies no original sample or estimate
here. Its normative point is that ambiguity-sensitive preferences might
challenge a proposed standard rather than merely demand a psychological
explanation of irrationality. The original Ellsberg treatment and
finite ambiguity models give their own versions and assumptions
(\cref{thm:p17-ellsberg-three-colors,def:p17-finite-maxmin});
the introduction's white/black labeling is retained locally.

The authors compare normative theorizing with theory construction,
without calling normative intuitions experimental observations. Initial
assessments, proposed rules and general principles can conflict; any
of them may require revision in a \emph{wide reflective equilibrium}.
Their swan, discarded-ticket and political-consistency examples are
starting judgments whose force depends on circumstances, not
incontrovertible axioms. A merely narrow equilibrium that fixes all
initial judgments cannot expose their possible defects. Their defense
of wide equilibrium invokes continuing argument and does not supply
an algorithm guaranteed to converge or a proof that every equilibrium
is uniquely correct.

\subsection{Descriptive mechanisms and their evidential standing}

The source organizes descriptive accounts around competence, departures
and interaction with prior beliefs. Rule-based accounts posit encoded
premises and inference schemata coordinated in working memory.
Errors may arise in interpretation, rule selection, coordination or
execution. Probabilistic computational accounts can describe an
input--output relation without specifying the cognitive operations
that implement it. Mental-model accounts instead posit representations
of possible situations and the construction and checking of alternatives;
missed alternatives and working-memory limitations offer different
error explanations. These are different proposed mechanisms, not
different names for an observed answer
\cite[pp.~21--23]{KnauffSpohn2021FrameworksIntro}.

The authors report supporting studies and implementations as a
literature-level overview. No common participant sample, experimental
protocol or discriminating model comparison is supplied in this
introduction. Its dual-process discussion likewise contrasts an
influential fast/slow account with empirical and theoretical objections
and later qualifications. It does not warrant identifying every fast
response with irrationality or every effortful response with correctness.
Prior beliefs enter all the main accounts, though their assigned roles
differ: they can moderate an otherwise abstract inference, or supply
much of the information retrieved to address it. An observed content
effect alone does not decide between those mechanisms.

For practical judgment the introduction describes Brunswik's lens
model: people select and weight cues whose relations to the environment
affect judgment. Its ecological-validity comparison concerns cues in a
specified environment, not an unconditional virtue of any highly
correlated cue. Judging an alternative's value differs from deciding
on an intended course of action. Prospect theory, bounded rationality
and other accounts address departures from classical economic models,
with controversy over which departures should count as errors. The
introduction also excludes its examples of bare perceptual selections
from the central assessment of intentional, reason-responsive choice;
that scope decision does not deny their scientific interest
\cite[pp.~23--24]{KnauffSpohn2021FrameworksIntro}.

\subsection{Autonomy and defeasible comparison}

The introduction rejects both making the two inquiries wholly
independent and treating their common word \emph{rationality} as
mere equivocation. It also resists naturalistic reductions. A predefined
success criterion can support an empirical inquiry into effective
means, but choosing that criterion as the proper aim remains a
normative question. The authors apply that objection to their readings
of Quine, Schurz and Hertwig, and Cohen. They do not deny that
empirical facts can matter to normative deliberation
\cite[pp.~24--26 and notes~28--33]{KnauffSpohn2021FrameworksIntro}.

Their proposed relation is \emph{defeasible}, not deductive. Human
receptivity to endorsed norms gives one possible reason to expect
conformity; observed considered judgments can give one possible reason
to question a proposed norm. Neither direction is automatic. Lack of
knowledge, conflicting motives, limited capacity, a mistaken standard
or a misdescribed task can defeat the inference. The source gives no
numerical reliability for this bridge and expressly acknowledges
counterexamples.

When evidence and benchmark differ, the authors identify three responses:
retain the norm and explain the departure; change the applicable
benchmark while retaining each benchmark's standing for its own task;
or treat the discrepancy as criticism of the normative theory itself.
The source's Wason/Bayesian and Popper--Kuhn/Ellsberg examples illustrate
these different moves. A smaller mismatch is not by itself proof that
the new norm is justified. Its \emph{double reflective equilibrium}
requires descriptive and normative theories to be considered together
without reducing one to the other; it leaves the arbitration of
competing norms genuinely difficult.

\section{Individual and social rationality}
\label{sec:p11-ri-individual-social}

The individual/social distinction cuts across both theoretical/practical
and normative/descriptive distinctions. One can assess an individual's
reasoning in a social setting, a group's formation of an opinion,
a jointly chosen action, or the organization of deliberation. These
are different targets. The authors acknowledge their introduction's
individualistic emphasis and describe social coverage as exemplary
rather than exhaustive. Note~35 preserves the distinction between
epistemology \emph{in} groups and epistemology \emph{of} groups
\cite[pp.~26--31 and notes~34--36]{KnauffSpohn2021FrameworksIntro}.

\subsection{Individual choice in a social setting}

Noncooperative game theory does not mean that players are forbidden
to cooperate. In the source's comparison, communication, bargaining
or agreement must be represented as actions when they matter to the
game. The cooperative approach instead starts with available
communication, possibilities of agreement and coalition formation,
with enforcement depending on the model; it is not a prediction that
every player will cooperate. Both labels concern a modeling approach rather
than a character judgment.

Nash equilibrium holds other players' strategies fixed and excludes
a profitable unilateral deviation
(\cref{def:p17-nash-equilibrium}). It does not optimize a collective
welfare criterion. The introduction also distinguishes classical,
Bayesian/epistemic and evolutionary interpretations rather than giving
one unconditional normative status to all game-theoretic models.
Utility need not encode selfishness: its content is a supplied
preference representation. An experiment rewarding only an individual's
money is not evidence that every social motive is self-regarding.

The source's one-shot prisoners' dilemma has mutual-cooperation payoff
\((2,2)\), mutual-defection payoff \((1,1)\), and unilateral
defection giving the defector 3 and the cooperator 0. Defection is
strictly dominant but both prefer joint cooperation to joint defection.
The earlier stipulated example in
\cref{def:p11-ri-two-agent-model,thm:p11-ri-individual-joint}
uses different cardinal payoffs with this same strict-dominance and
joint-improvement pattern; it is not a transcription of the source's
table or a claim that the cardinal models coincide.

Avoiding a Pareto-dominated outcome is also too weak to specify an
entire collective norm. In the source's table, each unilateral-defection
profile is Pareto efficient: improving the disadvantaged player would
lower the advantaged player's payoff. That does not establish fairness,
consensus or a just distribution. The authors mention Sen's liberal
paradox to indicate that even the necessity of Pareto efficiency is
contested in richer normative settings. The once-played dominance
calculation does not determine repeated-play incentives, enforcement,
communication or the interpretation of ultimatum-game responses.

\subsection{Joint belief, preference and action}

Epistemic social rationality concerns forming jointly justified beliefs,
conclusions or explanations. Practical social rationality concerns
joint decisions and action. The authors' examples include juries,
scientific collaboration, ethics committees and democratic institutions.
Their relevance does not establish that every such institution has
one mind or that individual belief-revision axioms transfer unchanged
to a collective. One must identify whether the group product is a
belief attribution, an accepted public commitment or an aggregation
of members' judgments.

Even coherent individual preference rankings need not produce a
coherent majority ranking. The following stipulated example makes
that limitation precise without importing a full impossibility theorem.

\begin{hbdefinition}{def:p11-ri-majority-model}{Three coherent rankings and pairwise majority}
Three voters strictly rank alternatives \(a,b,c\) as follows:
\[
 \begin{array}{c|c}
 \text{voter}&\text{ranking}\\ \hline
 1&a\succ_1 b\succ_1 c\\
 2&b\succ_2 c\succ_2 a\\
 3&c\succ_3 a\succ_3 b
 \end{array}
\]
Extend each row to its transitive strict ordering. Define
\(x\succ_M y\) when at least two voters rank \(x\) above \(y\).
These are ordinal comparisons; there is no common numerical welfare scale.
\end{hbdefinition}

\begin{hbtheorem}{thm:p11-ri-majority-cycle}{proved=annex:p11.ri.majority-cycle}{Individual transitivity does not ensure majority transitivity}
All three individual rankings are transitive, but pairwise majority gives
\(a\succ_M b\), \(b\succ_M c\) and \(c\succ_M a\).
Consequently majority fails transitivity in this profile.
\end{hbtheorem}

This counterexample concerns this aggregation rule and profile. It does
not prove that no aggregation procedure can work on any domain. The
introduction reports Arrow's much stronger impossibility under a
collection of conditions, not a prohibition on all collective decisions.
It also questions whether nondictatorship is a requirement of rationality
or a moral requirement. A formal incompatibility of desiderata cannot
choose which desideratum a community ought to relax.

\subsection{What the collective-reasoning evidence supports}

The source discusses demonstrability, or the capacity to exhibit a
correct solution to fellow group members, as one possible explanation
of successful epistemic collaboration. It distinguishes tasks announced
as having an objectively correct answer from tasks asking participants
to integrate opinions. Those different instructions can change the
social combination process. Its other communication approaches,
including Gricean, relevance-theoretic and argumentative accounts, are
reported alternatives rather than interchangeable mechanisms.

For one concrete example, the introduction cites Moshman and Geil's
1998 selection-task study. The original report identifies 143 students
from undergraduate educational-psychology courses, participating as one
way to fulfill a class requirement. Random assignment divided them into
32 individual solvers, 54 students working interactively in ten groups,
and 57 students first solving individually and then discussing in ten
groups. Each group had five or six members. These are three conditions,
not 143 independent group responses
\cite[pp.~233--234]{MoshmanGeil1998CollaborativeReasoning}.

The pictured cards were \texttt{E}, \texttt{K}, \texttt{4} and
\texttt{7}; the instruction was to establish whether the vowel/even
rule held for these four cards, with written justification. This has the
same countercase structure as \cref{thm:p11-ri-card-check}, but retains
the original study's different card labels. The experimenter observed
through a one-way mirror and videotaped the sessions. Groups were asked
to reach consensus; members then recorded and justified the answer
individually. The individual-then-interactive condition also recorded
each member's initial answer. At 25 minutes a group without consensus
was given ten more minutes; one unresolved group was then asked for
individual justifications and subsequently reached consensus. Individual
controls were encouraged to take sufficient time within their 50-minute
credit session; after their answer and justification they received the
correct solution and were asked to justify it
\cite[p.~234]{MoshmanGeil1998CollaborativeReasoning}.

\begin{hbconvention}{def:p11-ri-collaborative-statistics}{Response units in the original comparison}
In the reported tests below, \(N\) counts the stated response units,
\(\chi^2\) is the reported chi-square statistic, and \(p\) is its
reported significance probability, not a probability of a cognitive
mechanism or normative claim. A group decision counts once; its several
members' subsequent answers share the same discussion.
\end{hbconvention}

Exactly three of 32 individual controls selected only the vowel and
odd-number cards. Seven of ten interactive groups and eight of ten
individual-then-interactive groups selected that pattern. Thus the pooled
group success fraction was 15/20, or 75\%, against 3/32, or 9.375\%,
for controls. The authors report \(\chi^2(1,N=52)=23.42\),
\(p<.001\), for the pooled group/control comparison and
\(\chi^2(1,N=20)=0.27\), \(p>.05\), between the two group
conditions. The latter nonsignificance is not an equivalence result.
The measured product is a correct selection in this task, not a percentage
increase in a general rationality trait
\cite[p.~235, Table~1]{MoshmanGeil1998CollaborativeReasoning}.

Among the 57 individual-then-interactive students, 12 initially selected
the correct pattern and 45 finally did so. Of 37 changes in correctness,
35 were improvements and two deteriorations; the paper reports a binomial
\(p<.001\). These postdiscussion answers are not 57 independent
group-level trials. Seven of the ten groups contained an initially correct
member; five succeeded and two did not. The other three groups all
succeeded despite no member initially selecting the complete correct
pattern. In this sample, an initially correct member was therefore neither
necessary nor sufficient for success
\cite[pp.~236--237, Table~2]{MoshmanGeil1998CollaborativeReasoning}.

The authors interpret the advantage through joint construction of
arguments rather than simple copying of a complete answer. Their two
detailed discussion accounts show participants combining arguments about
the relevance of the odd card and the irrelevance of the even card;
neither group began with a correct complete selection. The third such
group's videotape was accidentally destroyed. The paper provides qualitative
sequences, not a formal coding protocol, inter-rater reliability estimate
or an experimentally isolated mediation effect
\cite[pp.~237--245]{MoshmanGeil1998CollaborativeReasoning}.

The original discussion considers time on task, majority conformity and
social influence as alternatives. Groups took approximately two and a
half times as long as individuals. The controls' available time and later
justification of the supplied answer inform that objection, but are not
a time-yoked manipulation. No initially correct majority or modal answer
existed in the ten individual-then-interactive groups; that rules out
mere adoption of such an initial vote in those groups, not every form
of social influence. The evidence supports the authors' task-specific
argument-construction interpretation with these limits. There was no
delayed test establishing durable individual learning, no transfer task
and no demonstration for other populations or arbitrary group sizes;
the authors explicitly decline the durable-general-change claim
\cite[pp.~242--246]{MoshmanGeil1998CollaborativeReasoning}.

More broadly, the introduction reports mixed group advantages.
Greater weight on a knowledgeable member, learning during interaction,
and aggregation of independent judgments are different candidate
explanations. A ``truth wins'' model supposes that the correct answer
can be recognized and adopted; it is not a theorem that the presence
of one knowledgeable member always suffices. Independent-crowd
aggregation assumptions do not survive simply by calling a
communicating group a crowd. Conformity, suppression of controversy,
selective sharing of common rather than unshared information, speaking
order and polarization can instead worsen the collective result.
These are the introduction's literature-level summaries, not a common
experimental dataset or a reconstructed causal account of the
Challenger disaster
\cite[pp.~29--31]{KnauffSpohn2021FrameworksIntro}.

\section{Process-oriented and output-oriented perspectives}
\label{sec:p11-ri-process-output}

An output-oriented account specifies which conclusion, epistemic state
or choice is appropriate given the input. A process-oriented account
specifies representations and operations mediating that relation.
The introduction relates this distinction to Marr's computational and
algorithmic levels, Newell's knowledge and symbol levels, and the
contrast between reason-based and mechanistic explanations. Those
comparisons are orientation, not proofs that the distinctions have
identical extensions in every theory. Mechanistic proposals identify
components and operations that can direct discovery and testing;
extending one to a new phenomenon requires checking both its relevant
similarities and its differences
\cite[pp.~31--34]{KnauffSpohn2021FrameworksIntro}.

The source also uses syntax and semantics as a historical illustration.
That illustration needs a precise logical boundary. First-order
soundness and completeness can identify derivability with semantic
consequence under the registered hypotheses
(\cref{thm:fol-completeness}). Incompleteness instead concerns
appropriate effectively axiomatized arithmetical theories and their
inability to decide every arithmetical sentence under the stated
consistency or soundness conditions
(\cref{thm:godel-first}). It does not show that every sound
proof system fails completeness for its own semantics. Equivalence of
the standard computation models supports the Church--Turing thesis
as an explication of effective calculation; it does not identify all
true propositions with algorithmically decidable ones.

Complex cognition involves interacting perceptual, linguistic, memory
and other processes. The introduction contrasts broad attempts to
model that interaction with more limited models of subcomponents and
the proposed computational-rationality program. Its historical account
does not establish that one program already explains all the involved
processes. It also distinguishes the internal evolution of thought
from externally prompted revision: choosing the appropriate new state
is an output question, but the route to that state can still require
a process explanation.

\subsection{Which conception of a state is being used?}

\begin{hbdefinition}{def:p11-ri-occurrent-dispositional}{Occurrent and dispositional attitudes}
In the source's philosophical usage, an occurrent attitude is presently
before the mind. A dispositional attitude is characterized through
responses it is disposed to produce under suitable elicitation, without
new information changing the underlying attitude. A disposition need
not be consciously considered at every moment.

An idealized deductively closed belief set or complete probability
measure need not be a list whose every item is simultaneously consciously
entertained. A modeler must say whether such an idealization specifies
dispositions, stored representations, an endorsed norm or currently
active contents. These interpretations are not made equivalent by
using the same mathematical object.
\end{hbdefinition}

The source tentatively connects occurrent attitudes with working-memory
representations and dispositional attitudes with long-term memory,
including global-workspace approaches and Goldman's discussion.
Note~37 identifies the philosophical background of dispositions.
The authors leave the match open: a philosophical disposition is not
automatically a particular memory trace, and a stored trace is not
automatically a presently endorsed belief. A check performed on active
contents does not certify consistency of all the agent's dispositions.
The memory account below develops the proposed cognitive interface
rather than converting that analogy into a universal identity.

\subsection{Two independent distinctions}

Neither perspective belongs exclusively to one discipline.
Psychological studies can compare inputs and choices without supplying
a process model. Philosophers and AI researchers can study possible
proof procedures, revision algorithms and decision-tree procedures
without making a psychological claim that people actually use them.
Moreover, the normative/descriptive distinction cuts across the
process/output distinction. The four resulting questions are:

\begin{enumerate}
\item Normative output: which conclusion or act should result under
      specified information, motives and an argued benchmark?
\item Descriptive output: which answers or choices actually result
      in the observed task and population?
\item Normative process: which procedures should be used, taking
      correctness, feasible resources and decision costs into account?
\item Descriptive process: which representations and operations
      actually produce the observed result?
\end{enumerate}

For a worked contrast, suppose two stipulated procedures return the
same correct answer, but the first costs one resource unit and the
second costs four. If only their answers are recorded, that observation
does not distinguish them. If the criterion is correct answer alone,
they tie; if the stipulated benefit is ten units and their known costs
are subtracted on a common scale, their net values are nine and six.
The latter comparison requires the cost model and evaluative criterion.
It is not participant data and does not show which procedure the
person used or that speed always outweighs other considerations.

The source invokes worst-case complexity as one possible normative
process consideration. Complexity is measured across a family of
encoded, finite instances of growing size, not an infinite menu within
one decision (\cref{def:p04-complexity-convention}). NP-hardness is
not an unconditional proof that every instance requires large effort;
a general polynomial-time solution to an NP-hard decision problem
under the stated reductions would imply \(\mathsf P=\mathsf{NP}\).
Small or specially structured instances can remain tractable.
Average-case claims additionally require an input distribution, and
approximation or satisficing changes the requested success criterion.
These qualifications preserve the introduction's point that translating
a formal worst-case result into a cognitive constraint requires more
than its complexity label.

Including decision costs in utility creates a further modeling issue:
determining the cost of choosing a procedure may itself require costly
deliberation. The introduction presents that potential circularity and
the satisficing response, not a theorem that all cost-aware decisions
are impossible. A supplied cost, as in the worked contrast, avoids
one dependency but does not explain how an actual agent obtains it.
The source therefore calls for normative process theories that address
how thinking should proceed, as well as whether its eventual answer
is satisfactory. The corresponding bounded-rationality distinctions
remain those of \cref{def:p17-br-principle-theory,def:p17-br-three-rationalities}.

\section{Preconditions of rationality}
\label{sec:p11-ri-preconditions}

The introduction asks about cortical, cognitive and social/evolutionary
conditions of human rationality. This is a biological and psychological
inquiry, not a proof that every possible rational system must have a
human brain. Multiple realizability and functionalist alternatives,
introduced earlier, are relevant precisely because implementation and
function are not identical questions
\cite[pp.~34--44]{KnauffSpohn2021FrameworksIntro}.

Its evolutionary illustration reports brain-volume ranges of about
400--500 cubic centimeters for \emph{Australopithecus africanus},
500--750 for \emph{Homo habilis} and 880--1,100 for
\emph{Homo erectus}; the latter is situated roughly 1.5--0.3 million
years ago. It cites a human worldwide average of about 1,350 cubic
centimeters and a broad healthy-adult range of approximately 900--2,000.
These are the cited survey's historical/anatomical summaries, not a
newly assembled dataset. Its sperm-whale, horse and elephant illustrations
concern absolute size, while its shrew illustration concerns size
relative to body mass. It also reports an approximately ten-percent
male/female mean volume contrast beside the absence of an average
traditional-intelligence-test difference. These are source summaries
without a reconstructed sampling or analysis record here. Their role
is to deny that absolute or body-relative size alone establishes
a general rationality ranking.

Likewise, the source's approximately ten-percent account of the
brain-volume/cognitive-ability association is not ten percent of a
person's rationality, nor a causal intervention effect. Its cited
work concerns intelligence differences, with further associations
involving regional anatomy, connectivity and energy use. The
introduction supplies no common sampling frame, operational definition
of rationality, estimate interval or adjustment set for those
comparisons. Its point is that volume alone is an inadequate explanation,
not that one anatomical measure settles every cognitive or social
assessment \cite[p.~35]{KnauffSpohn2021FrameworksIntro}.

\subsection{Cortical functions and their evidence}

For theoretical reasoning the authors organize their surveyed findings
around three broad networks. Language-associated left frontal and
prefrontal regions, including Broca's area, figure in the encoding and
processing of verbal material. Parietal, inferior temporal and occipital
regions figure in visual or spatial representation and imagery. Right
lateral/dorsolateral prefrontal regions figure in executive control and
conflict-related processing, including conflicts between a task's
formal validity and the plausibility of its content. This is their
organizing account, not an exhaustive one-region/one-function map
\cite[pp.~35--37]{KnauffSpohn2021FrameworksIntro}.

An activity contrast under a task and baseline is not by itself a
causal explanation or a demonstration that a region contains the
normatively correct answer. The introduction recognizes that limit
and discusses lesion comparisons and double dissociations as additional
evidence. A selective deficit must still be distinguished from
co-occurring damage, baseline differences, compensation and task
demands; a double dissociation alone does not identify two isolated
necessary-and-sufficient modules. Patient series may be small or even
single cases. Stimulation gives another intervention but is not
literally equivalent to deleting a brain region, and a change in
performance does not alone identify the representational format.
The source's own stimulation references include effects on reasoning,
not a uniform direction of impairment.

The introduction also reports a study by Gazzo Castañeda and colleagues
varying deductive versus probabilistic instructions, high/low conditional
probabilities and valid/invalid inference forms, while allowing binary
or graded answers. It reports different activity patterns and
content effects under those instructions, and binary responses even
among some probabilistically instructed participants. These reported
results motivate the possibility of task-dependent use of prior
knowledge rather than a universal graded-response preference. The
reference trail identifies the 2021 work as a \emph{manuscript submitted
for publication}, not an already qualified published experiment.
The introduction gives no sample, allocation, imaging-analysis or
uncertainty record sufficient to reconstruct that study here.

On the practical side, the source distinguishes cognitive-control
functions---switching, inhibition, error/conflict handling and working
memory---from value-related processing. It associates these respectively
with dorsolateral/anterior-cingulate and ventral/medial prefrontal
networks, while emphasizing interaction. Its lesion discussion
strengthens the functional hypothesis but does not establish that
every rational choice requires an identical localization. The actual
Goel treatment distinguishes observed activity from proposed inferential
function and preserves conflict-detection/resolution uncertainty
(\cref{sec:p19-rationality-1-4}). That interface does not automatically
qualify every study summarized in this introduction.

\subsection{Cognitive capacities, reflection and memory}

The source distinguishes rationality from intelligence as measured
through individual-difference tests. A level of cognitive capacity may
be required for a task without guaranteeing sound reasoning in it.
Its Stanovich--West illustration reports dependence on instructions,
including warnings to avoid biased processing. The task, warning,
measure of ability and rationality criterion therefore matter; a
correlation in one condition is not an identity of the constructs.
The source's discussion of ideological misuse of intelligence concepts
is historical caution, not a proof that every rationality concept is
politically neutral \cite[pp.~37--39]{KnauffSpohn2021FrameworksIntro}.

Metacognitive \emph{monitoring} assesses a cognitive activity, while
\emph{control} changes it or its resource allocation. Feeling that an
answer is right can influence reconsideration without guaranteeing
correctness. The introduction discusses implicit fluency/accessibility
cues, explicit assessments of one's abilities, confidence judgments
and the feeling of rightness. Its Thompson illustration obtains an
initial response, a feeling-of-rightness judgment and a later response;
that task sequence does not make longer reflection automatically
beneficial. The current meta-reasoning treatment separates cue validity,
cue use and the stage of judgment
(\cref{def:p19-r86-monitoring-control,def:p19-r86-judgment-types,
def:p19-r86-cue-standing}).

Philosophical higher-order attitudes supply a related output-oriented
inquiry. Believing something and believing that one believes it need
not be identified; requirements of self-awareness and second-order
infallibility remain disputed. Desires about desires and effective
second-order volitions also enter the source's Frankfurt comparison.
Auto-epistemic questions ask how present attitudes relate to beliefs
about future or past attitudes. Expecting to believe tomorrow that it
will rain does not, without reliability and information assumptions,
force belief in rain now: one might expect tomorrow's belief to be
based on misleading information. Interpersonal higher-order beliefs
and common knowledge further connect reflection with communication.
Neither the syntax of an iterated attitude expression nor a report
of confidence proves that the relevant attitude is rationally accurate.

Working memory maintains task-relevant information for active
processing; long-term memory supports stored knowledge and experience.
The source emphasizes capacity, interference, rehearsal and retrieval
conditions rather than one fixed duration. Its Miller illustration of
roughly seven elements is expressly followed by qualifications about
materials and task demands, including word length. It is not an
invariant seven-proposition limit on all reasoning. The source describes
decay, replacement, interference and consolidation as proposed memory
processes; failure of immediate recall does not by itself demonstrate
irretrievable loss or the absence of a stored disposition
\cite[pp.~39--40]{KnauffSpohn2021FrameworksIntro}.

Semantic memory concerns such knowledge and concepts as a pesticide's
classification; episodic memory includes personally situated events.
The source does not reconstruct every theory of those systems.
Its reasoning implication is that premises must be maintained, retrieved
and coordinated: losing one relevant premise can change the available
problem. Mental-model accounts emphasize the burden of constructing
alternatives; domain-knowledge accounts emphasize retrieved information.
Both need an account of the relevant representations and operations.
A consistency norm for an entire idealized belief set does not supply
unlimited access to every stored item. The proposed connection to
occurrent/dispositional attitudes remains a model-specific bridge,
not an identity deduced from memory terminology.

\subsection{Language, imagery and representation}

The source contrasts traditions identifying thought with inner speech
with accounts using mental models, spatial representation, simulation
or embodied experience. Its Sapir--Whorf discussion places linguistic
relativity and universalist responses beside one another, including
the Li--Gleitman and Levinson and colleagues spatial-reasoning dispute.
The described controversy is not settled by one's preferred nature/
nurture narrative or by a single language-associated activation.
Language can matter substantially without being necessary for every
form of reasoning in every child or nonhuman animal
\cite[pp.~40--42]{KnauffSpohn2021FrameworksIntro}.

Imagery adds a different representational question. Experiencing a
picture-like image is not evidence that every underlying operation
is pictorial. The introduction reports mixed effects of imaginable
content and distinguishes detailed visual imagery from more abstract
spatial or possibility representations. Its account allows irrelevant
visual detail to burden a task while a schematic model assists it.
A mental model of alternatives therefore need not resemble a photograph.
Introspection about the resulting experience does not disclose all
the processes that generated it.

The normative use of pictures is a further issue. A diagram can help
explain a relation without, solely by looking convincing, proving it.
Conversely, a diagrammatic proof system with specified syntax,
interpretation and licensed transformations is not excluded merely
because its expressions are pictures. The source identifies such
work through Jamnik's contribution while separating it from the
psychological question of how people experience imagery.

\subsection{Social and evolutionary hypotheses}

The introduction gives social coordination an important proposed
evolutionary role. Cosmides and Tooby's social-exchange account appeals
to checking reciprocal obligations and detecting cheaters, with Trivers's
reciprocal-altruism background. Its hunting/village illustration
explains the proposed function: people guarding the families and people
sharing the catch must track whether each side fulfills its commitment.
That explanatory scenario is not observed evidence about a particular
prehistoric contract or proof of a unique evolved module
\cite[pp.~42--44]{KnauffSpohn2021FrameworksIntro}.

Mercier and Sperber's argumentative account instead emphasizes
producing and assessing reasons in social exchange, and explains some
one-sided or unreliable solitary reasoning through that proposed
function. The introduction accepts sociality as important while
retaining the just-so-story objection: a plausible adaptive narrative
alone does not establish its actual history or exclude competitors.
Questions about language, metacognitive reflection and nonhuman
abilities remain both empirical and conceptual.

The source's customs-officer selection-task illustration states a
conditional linking an entry form with cholera listed on its other
side, in a protective-rule context. Its visible cases are entry,
no entry, cholera listed and cholera not listed. Under the intended
rule-enforcement interpretation, entry without the required list
item is the violating pattern, corresponding to the antecedent/
false-consequent check proved earlier
(\cref{thm:p11-ri-card-check}). The introduction reports improved
performance in that social context and a further supervisor/observer
role contrast. It does not supply original samples or uncertainty,
and those context effects do not by themselves establish a particular
evolutionary genealogy, universally eliminate errors, or show that
every conversational conditional is material implication.

The game-theoretic strand considers signaling, stag hunts and prisoners'
dilemmas. Cooperation can arise without being evolutionarily stable
in every model; an update rule and population/interaction structure
are needed for a stability claim. Tit for Tat can perpetuate retaliation,
whereas the source's GRIT illustration publicly signals willingness
to reconcile and offers bounded concessions in expectation of
reciprocation. That is a described conflict-reduction strategy, not a
guarantee against exploitation or a recommendation for every conflict.

Darwin's historical illustration in \emph{The Descent of Man},
volume~I, page~157, combines technological capacities with mutual aid
as responses to human physical limitations
\cite[p.~157]{Darwin1871DescentVolOne}.
This limited historical illustration does not establish every empirical
or social claim in that work. Neither natural selection
nor that passage entails a normative requirement of egoism.

\begin{hbexample}{def:p11-ri-false-belief}{The source's unexpected-transfer illustration}
Mary puts a candy in a red box and leaves. Another child moves it to
a blue box while an observing child watches. The observer is asked
where Mary will look. In the intended interpretation, Mary did not
observe or learn about the transfer, retains her original location
belief, wants the candy and searches according to that belief.
The predicted answer is therefore the red box although the candy is
actually in the blue one.

This separates the observer's information from the represented
information of another person. It is not an unconditional prediction
of Mary's behavior: changed knowledge, forgotten location, different
goals or another search policy can change it. No age-specific pass
rate or inference that every failing child lacks all understanding of
other minds is supplied by this illustration.
\cite[p.~43]{KnauffSpohn2021FrameworksIntro}
\end{hbexample}

The introduction connects such tasks with Tomasello's work on
intentional understanding and comparative cognition. Its proposed
priority of reflection on others' attitudes over reflection on one's
own is an evolutionary hypothesis, not established by the simple
box example. The account preserves social cognition as a candidate
precondition while declining to infer a unique human/nonhuman
rationality scale from a task or an animal-language assumption.

\section{Frontiers and neighboring inquiries}
\label{sec:p11-ri-frontiers}

The source asks whether domain-specific rationalities are necessary
and how rationality relates to neighboring concepts not systematically
treated in its own handbook. Its preference is for general frameworks
where their actual conditions apply, not the claim that every domain
has already been exhaustively explained. Generality still requires
the proper task, content, institutional conditions and normative aim
\cite[pp.~44--50]{KnauffSpohn2021FrameworksIntro}.

\subsection{Scientific, communicative and artificial-agent cases}

Scientific inquiry pursues justified beliefs, but selecting research
goals, funding priorities and beneficiaries also involves practical
values. The source contrasts vaccine research with cosmetic implants
to illustrate a value-laden choice, not to settle a research-allocation
policy by a scientific fact. Its Weber/Habermas value-freedom debate
therefore distinguishes epistemic methodology from the normative
organization and purposes of science
\cite[pp.~44--45 and note~38]{KnauffSpohn2021FrameworksIntro}.

On the authors' account, empirical findings support a normative
conclusion only with a further normative argument; their defeasible
bridge is not an exception granting deductive entailment without
premises. Transparent premises and public justification remain
important. Scientific expertise in hypothesis testing also does not
give every scientist authority over every practical value choice.
The source recognizes p-hacking and selective publication as possible
epistemic deficits rather than denying scientific competence wholesale.
Note~38 discusses the 2015 Open Science Collaboration and cautions
about interpreting nonreplication, including regression-to-the-mean
considerations and differences between subdisciplines. No replication
percentage, common failure cause or original-study reanalysis is
established in this introduction's evidence account.

Whether scientific inference uses irreducibly special norms remains
open in the source. Ockham's razor, statistical inference and inference
to the best explanation are candidate methodologies; belief revision
and Bayesian learning can offer general reconstructions. The historical
Popper--Kuhn/Lakatos dispute and the classical/Bayesian statistics
comparison exemplify competing proposals, not a proof that one
methodology exhausts science or everyday rationality.

Communicative rationality similarly has different conceptions.
Habermas's public-justification ideal involves nondomination and
unprejudiced consideration of arguments, exceeding mere effectiveness
at getting another person to act. Brandom's reason-giving practice
connects inferential relations with meaning; Grice's cooperative
principle and intentional account address how conversational
understanding works; Lewis's conventions supply a strategic account.
The introduction treats them as relationships to general epistemic
and practical rationality, not one identical doctrine
\cite[pp.~45--46]{KnauffSpohn2021FrameworksIntro}.

The authors also propose a Kantian reading of recognizing other
agents' aims in game theory. Recognition alone, however, does not
entail nondomination or moral respect: one can model another's aims
to manipulate that person. Equating a strategic representation with
Habermas's justificatory ideal therefore requires additional normative
conditions. This preserves the proposed comparison without claiming
that any equilibrium implements a discourse free of domination.

For AI, the introduction contrasts symbolic knowledge/procedure
representations with statistically trained pattern-learning systems.
It discusses explicit knowledge bases and revision, nonmonotonic
inference, expert-system limitations, predictive decision systems and
explainability as an interdisciplinary interface
\cite[pp.~46--47]{KnauffSpohn2021FrameworksIntro}.
Those are the authors' 2021 descriptions, not a benchmark of all later
AI systems. A symbolic representation does not guarantee correct
premises, and consistency is not a sufficient epistemic criterion.
The intended nonclassical logic and revision policy must be stated
rather than requiring every symbolic system to be classical.

Likewise, predictive success on a dataset neither establishes justified
belief in every new situation nor establishes a morally acceptable
decision criterion. The source's concern about opaque learned systems
is preserved as a demand for an account of the process and reasons;
its sweeping suggestion that no one can understand their decisions
is not adopted as a universal impossibility claim. An implementation
trace, statistical explanation and normative justification are different
objects. Hybrid symbolic/learning systems are one proposed route to
greater transparency, not a proof that all hybrids are rational or
that every explanation requires that architecture. No new AI-specific
criterion is established merely by classifying a system as intelligent.

\subsection{Emotion, morality and legal judgment}

The source separates three emotional-reasoning questions: effects of
the reasoner's current affect, effects of emotional task content and
their interaction. It reports contrasting findings from terror-related
content after the 7 July 2005 London attacks and from laboratory
mood/anxiety paradigms. These differ in exposure, content and conditions;
the introduction does not provide their original samples, effect sizes
or a common protocol. Ethical limits on inducing strong emotion and
differences between transient moods and persistent sentiments are
proposed explanations of disagreement, not established explanations
for each reported result. A ``cool head'' slogan cannot replace
that evidence \cite[pp.~47--48]{KnauffSpohn2021FrameworksIntro}.

Emotion may burden processing, guide attention to danger or contribute
to practical valuation. The source discusses arousal-and-interpretation
accounts while leaving the exact relationship between adequate emotions
and rationally justified emotions disputed. An emotion's having a
cause does not alone make it reason-responsive; conversely, calling
it reason-responsive does not settle every normative assessment of it.
Reports about impaired decision making in particular patients do not
prove that every emotion is helpful or that all emotion is irrational.

Moral sentimentalism, hedonic accounts and theories of happiness provide
different proposed relations between feeling, values and morality.
The source distinguishes ancient eudaimonia from a merely psychological
pleasant state, and contrasts Hume/Schopenhauer, Bentham and Hobbes
with Kant's noninstrumental account of practical reason. It does not
deduce a single moral theory from that comparison. The distinction
between a rational means and a justified ultimate end remains active.

Legal judgment provides one applied contrast. The introduction reports
Gazzo Castañeda and Knauff's work on legal conditionals, including
how mitigating or exculpatory conditions can conflict with an urge
to punish an act viewed as exceptionally immoral. This is a report
about particular tasks, not an independently qualified finding that
all professionals obey law and all laypeople disregard it. Which
legal rule applies, how a participant interprets the facts and whether
the rule itself is justified remain separate questions.

The trolley illustrations make the moral/empirical separation vivid.
In the stipulated switch case, intervention redirects a threat from
five people to one. In the footbridge variant, a person is used to
stop the trolley. Counting deaths yields the same five-versus-one
comparison, but intention, means, acts/omissions and personal involvement
can change the moral argument. The source's simplified contrast should
not be read as a theorem that every deontological account forbids
switching, or that every consequentialist comparison depends only
on fatalities \cite[pp.~48--49]{KnauffSpohn2021FrameworksIntro}.

Participants' reported choices and neural activity answer different
questions from which act is morally required. The source relates
Greene's imaging work to dual-process theories while retaining
objections to its interpretation: activity associated with emotion
does not establish irrationality, and activity associated with control
does not establish moral correctness. Its 2018 Moral Machine report
illustrates the danger of moving from online response frequencies
to a policy for autonomous systems without a normative argument.
Publication in 2018 is not asserted to be the platform's launch date;
the illustration is not a recommendation for vehicle or weapon design.

\subsection{Cultural, political and philosophical scope}

Plato's political organization, Hobbes's sovereign, Rousseau's social
contract and Rawls's veil of ignorance connect rational deliberation
with proposed accounts of political authority and justice. These
are selected source illustrations, not interchangeable definitions of
rationality or reconstructed proofs of their political conclusions.
The introduction deliberately does not claim a full political-philosophy
treatment \cite[pp.~49--50]{KnauffSpohn2021FrameworksIntro}.

Its intercultural discussion also marks a limitation of the selected
framework. The authors describe a literature search constrained by
languages they could read and acknowledge a predominantly European/
North-American intellectual orientation. Their references include
traditional Chinese yin/yang and the five phases---water, wood, fire,
metal and earth---culturally varying numerical systems,
Islamic economic ethics and a meditator/ultimatum-game comparison.
Those examples do not establish an invariant feature of every member
of a culture or a contradiction between number-system conventions
and mathematical truth. The authors' question about whether the
normative/descriptive dichotomy is equally explicit in every tradition
does not, solely through a different vocabulary, supply a deductive
bridge from fact to obligation. Different neural activity does not,
by itself, identify different normative standards.

The WEIRD-sampling concern is a generalization warning: responses
from Western, educated, industrialized, rich and democratic settings
cannot establish universality merely by being measured precisely.
Neither that warning nor the existence of cultural differences proves
unrestricted normative relativism. The same task, translation,
recruitment and inferential conditions must be considered before
making a cross-population comparison. The introduction itself
does not supply such an exhaustive comparison.

Finally, the source names neglected continental-philosophical
contributions: Hegel's dialectic, Nietzsche's perspectival critique,
Heidegger's criticism of technical rationality, Foucault's analysis
of the delimitation and government of reason, and Habermas's broader
democratic intentions. A dialectical use of conflicting claims is
not automatically a proof that an inconsistent classical premise
set avoids explosion. The analytic/continental labels also differ
in basis and hide internal variation.

The authors justify their editorial emphasis through the intersection
of formal analytic philosophy and cognitive psychology, and the
difficulty of experimentally addressing wider cultural/political
questions. That scope decision is not evidence that excluded traditions
lack sophisticated arguments or substantive contributions. The
boundaries of their 2021 volume therefore delimit this attributed
source treatment; they do not impose those exclusions on the whole
of \emph{The Methods of Reason}.

\section{Conclusion and open questions}
\label{sec:p11-ri-conclusion}

The introduction's conclusion is a research orientation, not a verdict
that humans are universally rational or universally irrational.
The theoretical/practical distinction identifies different attitudes
and justificatory roles. Normative/descriptive inquiry identifies
standards and observed performance without identifying either with
the other. Individual/social inquiry distinguishes agents in a shared
setting from collective products. Process/output inquiry distinguishes
what results from how it is realized, with both normative and
descriptive questions on either side
\cite[pp.~50--51]{KnauffSpohn2021FrameworksIntro}.

The authors' open agenda includes broader epistemic formats beyond
exclusive fixation on logic or probability, alternatives to standard
decision theory, genuinely collective rather than merely competitive
social accounts, and better process-oriented theories. They also
acknowledge unresolved accounts of concepts, meaning and truth.
Their double-equilibrium proposal asks the two disciplines to remain
responsive to one another while retaining normative argument and
empirical scrutiny as distinct activities. This is an advocated
methodology, not a convergence theorem or an experimental validation
of every proposed bridge.

The conclusion calls for contributions from linguistics, education,
anthropology and ethnology, including how rational thinking can be
developed and how its practices vary with context. It leaves the
initial rational-animal question deliberately unanswered. A diversity
of capacities and failures requires specified assessment objects,
information, aims and evidence; no single illustration answers it
for all humans or all other species.

\begin{hbmethod}{def:p11-ri-framework-comparison}{Use the framework without collapsing its distinctions}
For a reasoning problem, identify the assessed attitude or action and
whether the target is an individual in a social context or a collective
product. State the qualitative, comparative or quantitative format,
the current information and motives, and any revision input. Separate
the proposed norm and its justification from an empirical response
and an account of the process that produced it. Specify whether a
state is occurrent, dispositional or an idealized formal object, and
which cognitive, social and resource conditions are assumed.

Where observation and standard disagree, ask whether the task,
process explanation, benchmark application or norm itself is at issue;
give the evidence and argument for that move. For collective tasks,
state the aggregation or deliberation procedure and its success
criterion rather than inheriting an individual norm without a bridge.
For a specialized domain, retain the actual extra assumptions instead
of either inventing a new rationality solely from its name or forcing
it into an unsuitable general model. A finite contradiction or
counterexample establishes its specified boundary, not a population's
complete competence or an unconditional prescription.
\end{hbmethod}

\begin{hbwarning}{def:p11-ri-survey-evidence-boundary}{Source treatment and original evidence are different extents}
This chapter reconstructs the selected introduction's arguments,
examples, distinctions and acknowledged limits. Its empirical survey
reports remain attributed to that 2021 source unless a separate primary
evidence record is explicitly identified. No missing original sample,
task detail, analysis or uncertainty is inferred from a literature
summary, a title, a neural localization or an abstract.

In particular, the Moshman--Geil study distinguishes observed task
improvement from an isolated causal mechanism or durable transfer. The
introduction's submitted instruction/imaging study and other survey
reports require their original methods and analyses before their
findings can support stronger empirical conclusions. These evidence
boundaries prevent interpreting the
framework as a demonstrated universal cognitive model or prescribing
a disputed norm solely because some responses fit it.
\end{hbwarning}
